% BHU PYQ -- Philosophy: Ethics (2025-26 NEP)
\documentclass[11pt,a4paper]{article}
\usepackage[a4paper,top=2.5cm,bottom=2.5cm,left=2.8cm,right=2.8cm,headheight=14pt]{geometry}
\usepackage{amsmath,amssymb}


\usepackage{fontspec}
\setmainfont{Kohinoor Devanagari}
\usepackage{microtype,fancyhdr,enumitem,hyperref,lastpage,parskip}

\hypersetup{colorlinks=true,urlcolor=black,linkcolor=black,pdftitle={PHIMJ31: Ethics -- BHU NEP 2025-26}}

\pagestyle{fancy}\fancyhf{}
\renewcommand{\headrulewidth}{0.4pt}\renewcommand{\footrulewidth}{0.4pt}
\fancyhead[L]{\small \textit{Banaras Hindu University}}
\fancyhead[R]{\small \textit{PHIMJ31: Ethics (2025-26)}}
\fancyfoot[C]{\small Page \thepage\ of \pageref{LastPage}}

\newlist{parts}{enumerate}{2}
\setlist[parts,1]{label=(\alph*),leftmargin=*,itemsep=4pt,topsep=3pt}
\setlist[parts,2]{label=(\roman*),leftmargin=*,itemsep=2pt,topsep=2pt}

\newcommand{\pts}[1]{\hfill \textit{[#1]}}

\begin{document}

\begin{center}
  {\large\bfseries BANARAS HINDU UNIVERSITY}\\[2pt]
  {\normalsize B. A. (Hons.) Semester III Examination, 2025-26 (NEP)}\\[6pt]
  \rule{0.8\linewidth}{0.5pt}\\[6pt]
  {\large\bfseries PHILOSOPHY}\\[4pt]
  {\large\bfseries Paper Code: PHIMJ31 : Ethics}\\[4pt]
  {\normalsize Time: 2:30 Hours \hfill Full Marks: 70}\\[2pt]
  {\small REF NO. 2025-26/D-III/090}
\end{center}
\rule{\linewidth}{0.4pt}
\smallskip

\noindent \textit{Instructions: Attempt all the three Sections A, B and C according to instructions given. Write your Roll No.\ at the top immediately on receipt of this question paper.}

\smallskip
\noindent \textit{दिये गये निर्देशों के अनुसार खण्ड अ, ब तथा स के उत्तर दीजिये।}

\bigskip

\section*{SECTION -- A / खण्ड -- अ}
\noindent \textit{Note: Attempt any \textbf{three} questions out of five, each in 300 words. Each question carries 10 marks.}\pts{3 $\times$ 10 = 30}

\noindent \textit{निम्नलिखित पाँच प्रश्नों में से किन्हीं तीन प्रश्नों का उत्तर दीजिये, प्रत्येक उत्तर 300 शब्दों में हो। प्रत्येक प्रश्न 10 अंकों का है।}

\begin{enumerate}[label=\textbf{\arabic*.},leftmargin=*,itemsep=12pt]

  \item Discuss the nature and scope of Ethics.\pts{10}
  
  नीतिशास्त्र के स्वरूप तथा क्षेत्र की विवेचना कीजिये।

  \item Examine J. S. Mill's qualitative utilitarianism.\pts{10}
  
  जे० एस० मिल के गुणात्मक उपयोगितावाद का परीक्षण कीजिये।

  \item What is categorical imperative? Explain according to Kant.\pts{10}
  
  निरपेक्ष आदेश क्या है? काण्ट के अनुसार व्याख्या कीजिये।

  \item Give an account of the concept of `rta' as a moral principle.\pts{10}
  
  एक नैतिक सिद्धान्त के रूप में `ऋत' की अवधारणा का विवरण दीजिये।

  \item Discuss Dr. B. R. Ambedkar's critique of Hindu Ethics.\pts{10}
  
  डॉ० बी० आर० अम्बेडकर द्वारा हिन्दू नीतिशास्त्र की समीक्षा की चर्चा कीजिये।

\end{enumerate}

\bigskip

\section*{SECTION -- B / खण्ड -- ब}
\noindent \textit{Note: Attempt any \textbf{four} questions out of the following six questions, each in 150 words. Each question carries 5 marks.}\pts{4 $\times$ 5 = 20}

\noindent \textit{निम्नलिखित छह प्रश्नों में से किन्हीं चार प्रश्नों का उत्तर दीजिये, प्रत्येक उत्तर 150 शब्दों में हो। प्रत्येक प्रश्न 5 अंकों का है।}

\begin{enumerate}[label=\textbf{\arabic*.},leftmargin=*,start=6,itemsep=10pt]

  \item Discuss the idea of the good as a moral concept.\pts{5}
  
  एक नैतिक अवधारणा के रूप में शुभ के विचार की चर्चा कीजिये।

  \item Explain Aristotle's doctrine of the golden mean.\pts{5}
  
  अरस्तू के स्वर्णिम मध्य मार्ग के सिद्धान्त की व्याख्या कीजिए।

  \item Discuss the concept of freedom of will as a postulate of morality.\pts{5}
  
  एक नैतिक पूर्वमान्यता के रूप में संकल्प की स्वतन्त्रता की अवधारणा की चर्चा कीजिए।

  \item Explain Henry Sidgwick's rationalist utilitarianism.\pts{5}
  
  हेनरी सिजविक के बुद्धिवादी-उपयोगितावाद की व्याख्या कीजिए।

  \item What is retributive theory of punishment? Explain.\pts{5}
  
  दण्ड का प्रतीकारात्मक सिद्धान्त क्या है? व्याख्या कीजिए।

  \item Discuss Plato's concept of justice as virtue of the soul.\pts{5}
  
  प्लेटो के दर्शन में आत्मा के गुण के रूप में न्याय की अवधारणा की चर्चा कीजिए।

\end{enumerate}

\bigskip

\section*{SECTION -- C / खण्ड -- स}
\noindent \textit{Note: Attempt all the \textbf{ten} questions of this Section in a maximum of two sentences each. Each question carries 2 marks.}\pts{10 $\times$ 2 = 20}

\noindent \textit{इस खण्ड के सभी दस प्रश्नों के उत्तर अधिकतम दो वाक्यों में दीजिये। प्रत्येक प्रश्न 2 अंकों का है।}

\begin{enumerate}[label=\textbf{\arabic*.},leftmargin=*,start=12,itemsep=8pt]
  \item 
  \begin{parts}
    \item What does `right' signify as a moral category?\pts{2}
    
    एक नैतिक संप्रत्यय के रूप में `उचित' का क्या अर्थ है?

    \item What do you mean by Virtue ethics?\pts{2}
    
    सद्गुण नीतिशास्त्र से आप क्या समझते हैं?

    \item What is deontology?\pts{2}
    
    कर्तव्यशास्त्रीय नीतिशास्त्र क्या है?

    \item What do you understand by a teleological theory?\pts{2}
    
    प्रयोजनवादी नीतिशास्त्र से आप क्या समझते हैं?

    \item What is the ultimate end of human life according to Hindu ethics?\pts{2}
    
    हिन्दू नीतिशास्त्र के अनुसार मानव जीवन का परम लक्ष्य क्या है?

    \item What is the deterrent theory of punishment?\pts{2}
    
    दण्ड का निवारक सिद्धान्त क्या है?

    \item What is the criterion of universalisability?\pts{2}
    
    सर्वव्यापीकरण की कसौटी से क्या अभिप्राय है?

    \item What is hedonistic calculus?\pts{2}
    
    सुखवादी कलन क्या है?

    \item What do you understand by the principle of utility?\pts{2}
    
    उपयोगिता के नियम से आप क्या समझते हैं?

    \item What is the complete good according to Kant?\pts{2}
    
    काण्ट के अनुसार पूर्ण शुभ क्या है?
  \end{parts}
\end{enumerate}

\end{document}
